\documentclass[11pt]{scrartcl}

%%%%%%%%%%%%%%%%%%%%%%%%%%%%%%%%%%%%%%%%%%%%%%%%%%%%%%%%%%%%%%%%%%%%%%%%%%%%%%%%
% Make this compatible with pdflatex and lualatex, xelatex.
%%%%%%%%%%%%%%%%%%%%%%%%%%%%%%%%%%%%%%%%%%%%%%%%%%%%%%%%%%%%%%%%%%%%%%%%%%%%%%%%

\usepackage{ifxetex,ifluatex}
\newif\ifxetexorluatex
\ifxetex
  \xetexorluatextrue
\else
  \ifluatex
    \xetexorluatextrue
  \else
    \xetexorluatexfalse
  \fi
\fi

\ifxetexorluatex
  \usepackage{fontspec}
  \usepackage{xltxtra} % Not that important nowa
  % Font packages, and other xetex specific configuration here.
\else
  % Use utf8 as encoding.
  \usepackage[utf8]{inputenc}
  \usepackage[T1]{fontenc}
  % font packages
  \usepackage{lmodern} % Better than default computer modern?
\fi

%%%%%%%%%%%%%%%%%%%%%%%%%%%%%%%%%%%%%%%%%%%%%%%%%%%%%%%%%%%%%%%%%%%%%%%%%%%%%%%%
% End of pdflatex/xelatex/lualatex specific configuration
%%%%%%%%%%%%%%%%%%%%%%%%%%%%%%%%%%%%%%%%%%%%%%%%%%%%%%%%%%%%%%%%%%%%%%%%%%%%%%%%

\usepackage[hmargin=2cm,vmargin=2.5cm]{geometry} % Set small margins for handout

% BibTeX setup
\usepackage[backend=bibtex, bibstyle=alphabetic, citestyle=alphabetic]{biblatex}
\bibliography{references}

% \usepackage[english]{babel} % English
\usepackage[ngerman]{babel} % German

\usepackage{enumitem} % Package for handling of lists
\setlist{noitemsep}   % No separating whitespace between list items.

% Standard packages for math-related things.
\usepackage{amsmath}
\usepackage{amssymb}
\usepackage{amsthm}

\usepackage{mathtools}

\usepackage{graphicx} % to include graphics with \includegraphics[options]{imagefile}

\usepackage{tikz}
\usepackage{pgfplots}

% Hyperref is great, but sometimes there are problems if packages are loaded before or after hyperref. Check the documentation.
\usepackage{hyperref} % For PDF links, toc, etc.

% Formatting of paragraphs
\parindent 0cm                     % No intendation at the beginning of a paragraph
\parskip1.5ex plus0.5ex minus0.5ex % Vertical space between two paragraphs


%%%%%%%%%%%%%%%%%%%%%%%%%%%%%%%%%%%%%%%%%%%%%%%%%%%%%%%%%%%%%%%%%%%%%%%%%%%%%%%%
% Theorem and style definitions (for amsthm). For details see
% https://en.wikibooks.org/wiki/LaTeX/Theorems
%%%%%%%%%%%%%%%%%%%%%%%%%%%%%%%%%%%%%%%%%%%%%%%%%%%%%%%%%%%%%%%%%%%%%%%%%%%%%%%%

\theoremstyle{plain} % Usual style for theorems, etc.

% All numbered with the same counter (theorem).
% \newtheorem{theorem}{Theorem}[section] % This uses numbering SECTION.COUNT instead of COUNT. Useful in longer documents.
\newtheorem{theorem}{Satz} % Number linearly (no SECTION prefix).
% Usage: \newtheorem{environmentname}[counter]{displayedtext}
\newtheorem{proposition}[theorem]{Proposition}
\newtheorem{lemma}[theorem]{Lemma}
\newtheorem{corollary}[theorem]{Korollar}
\newtheorem*{lemma*}{Lemma} % not numbered.

\theoremstyle{definition} % Usual style definitions.
\newtheorem{definition}[theorem]{Definition}
\newtheorem*{definition*}{Definition}

\theoremstyle{remark} % Usual style definitions.
\newtheorem{remark}[theorem]{Remark}


%%%%%%%%%%%%%%%%%%%%%%%%%%%%%%%%%%%%%%%%%%%%%%%%%%%%%%%%%%%%%%%%%%%%%%%%%%%%%%%%
% Some commands
%%%%%%%%%%%%%%%%%%%%%%%%%%%%%%%%%%%%%%%%%%%%%%%%%%%%%%%%%%%%%%%%%%%%%%%%%%%%%%%%

\newcommand{\IN}{\mathbb{N}} % blackboard bold N for natural numbers
\newcommand{\IR}{\mathbb{R}} % blackboard bold R for real numbers
\newcommand{\IZ}{\mathbb{Z}} % blackboard bold Z for integers
\newcommand{\IF}{\mathbb{F}} % blackboard bold F for fields
\newcommand{\IG}{\mathbb{G}} % blackboard bold G for groups

% feel free to add more commands here

\newcommand{\set}[1]{\{#1\}}

%%%%%%%%%%%%%%%%%%%%%%%%%%%%%%%%%%%%%%%%%%%%%%%%%%%%%%%%%%%%%%%%%%%%%%%%%%%%%%%%
% Content begins here
%%%%%%%%%%%%%%%%%%%%%%%%%%%%%%%%%%%%%%%%%%%%%%%%%%%%%%%%%%%%%%%%%%%%%%%%%%%%%%%%

\def\Semester{WS 21/22}
%\def\Semester{SS 2021}
\def\Seminar{Mathematische Aspekte des maschinellen Lernens}
\def\Title{Approximationstheorie} 
\def\Author{Pavel Zwerschke} 
\def\Date{19.11.2021} % The date of your presentation

% Font
\renewcommand{\familydefault}{\sfdefault}


% Start of document
\begin{document}

\vspace*{-2cm}
\begin{minipage}{3cm}
%	\includegraphics[width=3cm]{./logos/kit-en.pdf} % English logo
	\includegraphics[width=3cm]{./logos/kit-de.pdf} % German logo
\end{minipage}\hspace*{0.2cm}~
\begin{minipage}{14cm}{
		\sffamily \Large{\Seminar} \hfill \Date \\ 
		\Author \hfill \Semester
}\end{minipage}
\vspace{-0.5cm}
\begin{center}
	\huge \sffamily \Title
\end{center}
\vspace{-1cm}

\section{A section}
\label{sec:a-section}

Some text in a section.

\section*{An unnumbered section}

This section, unlike section~\ref{sec:a-section} is not numbered.
This is indicated by the ``*'', i.e., \verb|\section*| instead of \verb|\section|.
Such ``starred'' and ``unstarred'' variants also exist for many environments (and commands).
There are subsections, subsubsections and paragraphs (and subparagraphs) for more fine-grained subdivision.

\subsection*{A subsection}

Some text.

\subsubsection*{A subsubsection}

More text.

\paragraph*{A paragraph}

Even more text.


\section{Important hint}

If you don't know \LaTeX{}, how to use math mode, environment, and so on,
a short tutorial/introduction will be more useful than this document.
This document is only a reminder of some basic \LaTeX{} stuff.
Not the absolute basics, and nothing advanced.

Furthermore, use a \textbf{suitable editor}.
This will make using \LaTeX{} \emph{much} easier.

\section{Basic usage of \LaTeX}
\label{sec:math-and-env}


\subsection{{\texttt{\textbackslash{}newcommand} for shorthands}}
\label{subsec:newcommand}

Using \verb|\newcommand{\cmdname}{output}| you can define shorthands,
e.g., \verb|\IN| for natural numbers $\IN$, and so on.
See the preamble.
Much more than mere text insertion is possible.
See (advanced) \LaTeX{} introductions.

\subsection{Lists}
\label{subsec:lists}

Use \verb|itemize|, \verb|enumerate| or \verb|description| for lists.
For example (itemize):
\begin{itemize}
	\item First item
	\item Second item
\end{itemize}


\subsection{Math environments}
\label{subsec:math-env}

Use \emph{inline} math mode (i.e., tex code in \$'s) for inline math,
e.g., $\lim_{x \to 1} e^{2 \pi i x} = 1$ or $\alpha \beta = \Gamma$.
Use the \verb|equation*| environment (or \verb|\[| and \verb|\]|)
for math in \emph{display} mode, e.g.,
\begin{equation*}
\lim_{x \to 1} e^{2 \pi i x} = 1.
\end{equation*}
If you use \verb|equation| you get numbered equations which you can reference to (see section~\ref{sec:references}).
\begin{equation}
\label{eq:eqn}
|x| = \left\{
\begin{array}{rl}
-x &\mbox{ if $x<0$} \\
x &\mbox{ otherwise}
\end{array}
\right.
\end{equation}

Proofs, Definition, Lemmata, Propositions, Theorems, Remarks, etc, have their own environments.
Environments can be be changed/redefined and new ones can be defined.

\begin{definition}[Negligible function]
	\label{def:negl}
	Let $f \colon \IN \to \IR$ be a function.
	If for any $k > 0$, it holds that $f(x) x^k \rightarrow 0$ for $x \to \infty$,
	then we call $f$ \emph{negligible}.
\end{definition}

\begin{lemma}
	$f(x) = 0$ is a negligible function.
\end{lemma}

\begin{proposition}
	If $f$ and $g$ are negligible functions,
	then $f + g$ is a negligible function.
\end{proposition}

The following theorem has a proof included.

\begin{theorem}[The ring of negligible functions]
	The set of negligible functions is closed under addition, subtraction and multiplication.
\end{theorem}
\begin{proof}
	This is just an application of theorems about limits of series and induction.
	For example, $0 = 0 + 0 = \lim_{x \to \infty} f(x)x^k + \lim_{x \to \infty} g(x)x^k = \lim_{x \to \infty} (f(x) + g(x))x^k$
	and
	\begin{equation*}
	0 = 0 \cdot 0 = \lim_{x \to \infty} f(x) x^{k} \lim_{x \to \infty} f(x) g(x)) x^{k} = \lim_{x \to \infty} (f(x) g(x)) x^{2k}
	\end{equation*}
\end{proof}

\begin{corollary}
	Working with negligible functions is easy.
\end{corollary}

\begin{remark}
	If $f$ is \emph{not} negligible,
	this does \emph{not} imply that $|f(x)| \geq x^{-k}$ always.
	This needs to hold for infinitely many $x \in \IN$ (and $k > 0$).
	E.g., $f(x) = 1 - (-1)^x$ is not negligible.
\end{remark}

There are a lot more useful things for math layout,
e.g., \verb|align| and \verb|aligned| environments for equations, and so on.

\subsection{References}
\label{sec:references}

To refer to a section, or any other ``referrable'' object, use the ``ref'' command.
For example: Section~\ref{sec:a-section} or Definition~\ref{def:negl} or Eq.~(\ref{eq:eqn}).
The tilde \verb|~| is an ``unbreakable space''.
(There are more advanced ways to do this, which are especially useful for longer documents.
For example the ``cleveref'' package.)


\subsection{Literature and BibTeX}
\label{subsec:literature}

To refer to literature, use BibTeX (which needs a bib-file) and run it (on the main document).
The bibfile is set by \verb|\bibliography{FILENAME}| in the preamble.
The literature inserted by the \verb|\printbibliography{}| command. (See the end of the document).
A good source of (sample) bibfiles is \url{https://dblp.uni-trier.de}.
Use \verb|\cite{}| (e.g., \cite{DH76} or \cite{RSA78}) or variants (see Section~\ref{subsec:packages}).


\subsection{Packages and the internet}
\label{subsec:packages}

There is a huge supply of useful packages.
For almost every problem, there's a package to solve it.
Just use the internet to find them.
(Special mention: TikZ, cleveref, cryptocode)

Reading some short introduction/tutorial on \LaTeX{} is also recommended.
Because this document only scratches the surface:
It does not have tables, pictures, splitting the document into multiple files, and so on.
Good starting points are:
Search engines, \url{https://en.wikibooks.org/wiki/LaTeX}, \url{https://tex.stackexchange.com}



\subsection{Miscellaneous}
\label{subsec:misc}

To start a new paragraph, use either an empty line in the source tex file or the command \verb|\par|.

Footnotes work via the \verb|footnote| command.\footnote{This is a footnote.}


%%%%%%%%%%%%%%%%%%%%%%%%%%%%%%%%%%%%%%%%%%%%%%%%%%%%%%%%%%%%%%%%%%%%%%%%%%%%%%%%
% Content ends here
%%%%%%%%%%%%%%%%%%%%%%%%%%%%%%%%%%%%%%%%%%%%%%%%%%%%%%%%%%%%%%%%%%%%%%%%%%%%%%%%
\printbibliography{}
\end{document}